% ================================================================
% ==== FILE: content/k_levels/k6.tex
% ================================================================

% ==============================
%  Ontology of Continua — Core
%  K-level Module: K6
%  FULL MODULE — FINAL
% ==============================

\subsubsection{\texorpdfstring{$K_6$}{K_6} Übersicht}
\label{sec:k6-overview}

$K_6$ ist die Ebene, auf der \emph{kognitive Organisation} entsteht
ontologische Dimension. Die erregbare elektrische Dynamik von $K_5$
stabile Attraktoren produzieren; Diese Attraktoren werden zu \emph{Modellen} dafür
kodieren interne Struktur, Gedächtnis, Vorhersagezustände und Repräsentation
Beziehungen.

Die Geburt von $K_6$ erfolgt, wenn die Anregungszyklen von $K_5$ nicht mehr vorhanden sind
lediglich elektrische Wellen, werden aber \emph{informativ interpretierbar}
Muster, deren Unterschiede nicht in $A(K_5)$ ausgedrückt werden können. Dies
erfordert einen neuen Satz von Achsen und eine neue Dimensionsstruktur.

Die minimalen Features:
\begin{itemize}
    \item repräsentative Zustände,
    \item-Vorhersagefähigkeit (Vorhersagefehlerachse),
    Semantische und Bindungsstruktur von \item,
    \item interne Modelle $\mu_t$,
    \item Speicherkonsolidierung/Rekonsolidierung,
    \item stabile Komponenten kognitiver Graphen,
    \item Entstehung von Protosymbolen.
\end{itemize}

$K_6$ ist das kleinste Kontinuum, das den Namen „kognitiv“ verdient. Es
ist noch nicht sprachlich oder sozial (diese gehören zu $K_7$).

% ================================================================
\subsubsection{Zustandsraum \texorpdfstring{$\Omega(K_6)$}{\Omega(K_6)}}

Ein Zustand von $K_6$ ist:
\[
\Omega(K_6) =
\{
m,\; \mu_t,\; G_{\mathrm{cog}},\;
x(t),\; r_i,\;
Haustier,\; B_t,\;
M_{\mathrm{con}}, M_{\mathrm{rec}}
\}.
\]

Elemente:
\begin{itemize}
    \item $m$ – aktives kognitives Modell,
    \item $\mu_t$ – Verteilung der Modellzustände,
    \item $G_{\mathrm{cog}}$ – kognitiver Graph (Knoten = Modelle, Kanten = Bindung),
    \item $x(t)$ – Aktivierungsmuster, geerbt von $K_5$-Attraktoren,
    \item $r_i$ – repräsentative Einheiten (Protosymbole),
    \item $PE_t$ – Vorhersagefehler,
    \item $B_t$ – Bindungskonfiguration,
    \item $M_{\mathrm{con}}$, $M_{\mathrm{rec}}$ – konsolidierte vs. rekonsolidierende Speicherzustände.
\end{itemize}

Zulässigkeit:
\[
D(m) \le \Theta_{\mathrm{cog}}, \qquad
|PE_t| < \Theta_{\mathrm{pred}}, \qquad
B_t \in \text{Bindungskapazität}.
\]

$D(m)$ ist eine Inkonsistenz des Modells $m$.

% ================================================================
\subsubsection{Grenze \texorpdfstring{$\partial\Omega(K_6)$}{\partial\Omega(K_6)}}

Zu den kognitiven Grenzen gehören:
\begin{itemize}
    \item Grenzen der Darstellungsfähigkeit,
    \item Kanten der Bindungstiefe,
    \item maximaler Vorhersagehorizont,
    \item Speicherstabilitätsgrenzen,
    \item semantische Kohärenzgrenzen,
    \item-Übergänge, bei denen $PE_t$ nicht abgeglichen werden kann.
\end{itemize}

Das Überqueren von $\partial \Omega(K_6)$ bricht die Modellierungsfähigkeit zusammen.

% ================================================================
\subsubsection{Achsen \texorpdfstring{$A(K_6)$}{A(K_6)}}

Wie in K6-Lauf 2 festgelegt, ist der Mindestsatz:
\[
A(K_6) = \{
A_{\mathrm{sel}},\;
A_{\mathrm{cmp}},\;
A_{\mathrm{bind}},\;
A_{\mathrm{cat}},\;
A_{\mathrm{mem}},\;
A_{\mathrm{pred}},\;
A_{\mathrm{Modell}}
\}.
\]

Interpretation:
\begin{itemize}
    \item $A_{\mathrm{sel}}$ — Auswahlachse (welches Modell aktiv ist),
    \item $A_{\mathrm{cmp}}$ — Vergleichsachse,
    \item $A_{\mathrm{bind}}$ – Bindungsachse (Merkmalszusammensetzung),
    \item $A_{\mathrm{cat}}$ – Kategorisierungsachse,
    \item $A_{\mathrm{mem}}$ – Speicherachse,
    \item $A_{\mathrm{pred}}$ – Vorhersageachse,
    \item $A_{\mathrm{model}}$ – Modellidentitätsachse.
\end{itemize}

Diese Achsen können nicht auf Erregungsunterschiede von $K_5$ reduziert werden.

% ================================================================
\subsubsection{Potentiale \texorpdfstring{$P(K_6)$}{P(K_6)}}

Potenziale regulieren Erkenntnisartefakte:

\[
P(K_6)=
\{
P_{\mathrm{pred}},\;
P_{\mathrm{val}},\;
P_{\mathrm{bind}},\;
P_{\mathrm{sal}},\;
P_{\mathrm{mem}},\;
P_{\mathrm{elim}}
\}.
\]

Wo:
\begin{itemize}
    \item $P_{\mathrm{pred}}$ – Vorhersagepotenzial zur Steuerung der Modellaktualisierung,
    \item $P_{\mathrm{val}}$ — Valenzpotential (affektive Gestaltung),
    \item $P_{\mathrm{bind}}$ – Potenzial für Feature-Bindung,
    \item $P_{\mathrm{sal}}$ – Salienzmodulation,
    \item $P_{\mathrm{mem}}$ – Konsolidierungsdruck für das Gedächtnis,
    \item $P_{\mathrm{elim}}$ – Eliminierung inkohärenter Modelle.
\end{itemize}

% ================================================================
\subsubsection{Schwellenwerte \texorpdfstring{$\Theta(K_6)$}{\Theta(K_6)}}

Schwellenwerte bestimmen die kognitive Lebensfähigkeit:
\[
\Theta(K_6)=
\{
\Theta_{\mathrm{cog}},\;
\Theta_{\mathrm{pred}},\;
\Theta_{\mathrm{bind}},\;
\Theta_{\mathrm{sal}},\;
\Theta_{\mathrm{Modell}},\;
\Theta_{\mathrm{mem-stab}},\;
\Theta_{\mathrm{noise-cog}}
\}.
\]

Bedeutung:
\begin{itemize}
    \item $\Theta_{\mathrm{cog}}$: globale Inkonsistenztoleranz,
    \item $\Theta_{\mathrm{pred}}$: maximal zulässiger Vorhersagefehler,
    \item $\Theta_{\mathrm{bind}}$: Bindungskapazitätsgrenze,
    \item $\Theta_{\mathrm{sal}}$: Sättigungsschwelle für Salienz,
    \item $\Theta_{\mathrm{model}}$: Kollapsgrenze für die Lebensfähigkeit des Modells,
    \item $\Theta_{\mathrm{mem-stab}}$: Speicherstabilitätsschwelle,
    \item $\Theta_{\mathrm{noise-cog}}$: kognitive Lärmtoleranz.
\end{itemize}

Das Überqueren dieser Werte führt je nach Richtung zum Zusammenbruch von $K_6$ oder zum Übergang zu $K_7$.

% ================================================================
\subsubsection{Flüsse \texorpdfstring{$J(K_6)$}{J(K_6)}}

Die Abläufe sind informativ und repräsentativ:

\begin{itemize}
    \item $J_{\mathrm{pred}} = -\nabla_m PE_t$ – Vorhersageaktualisierungsfluss,
    \item $J_{\mathrm{bind}}$ – Bindung/Entbindung von Übergängen,
    \item $J_{\mathrm{cmp}}$ — Vergleichsdynamik,
    \item $J_{\mathrm{sel}}$ – Modellauswahlfluss,
    \item $J_{\mathrm{mem}}$ – Konsolidierungs-/Rekonsolidierungsfluss,
    \item $J_{\mathrm{elim}}$ – Eliminierung inkohärenter Modelle.
\end{itemize}

Stabilität erfordert:
\[
\sigma(J_{\mathrm{pred}}) < \sigma_{\max}(K_6).
\]

% ================================================================
\subsubsection{Zyklen \texorpdfstring{$C(K_6)$}{C(K_6)}}

Zu den charakteristischen Zyklen gehören:

\[
C(K_6) =
\{
C_{\mathrm{pred}},\;
C_{\mathrm{bind}},\;
C_{\mathrm{mem}},\;
C_{\mathrm{cmp}},\;
C_{\mathrm{Modell}}
\}.
\]

Beschreibungen:
\begin{itemize}
    \item $C_{\mathrm{pred}}$: Vorhersagezyklus (Modell → Vorhersage → Fehler → Aktualisierung),
    \item $C_{\mathrm{bind}}$: Feature-Bindungszyklus (Zusammensetzung → Stabilisierung → Freigabe),
    \item $C_{\mathrm{mem}}$: Speicherzyklus (Konsolidierung → Speicherung → Rückkonsolidierung),
    \item $C_{\mathrm{cmp}}$: Vergleichs-Differenzierungs-Zyklus,
    \item $C_{\mathrm{model}}$: Zyklus der Modellauswahl, Bewertung, Eliminierung.
\end{itemize}

Jeder hat die charakteristische Periode $\tau_{C_j}$.

% ================================================================
\subsubsection{Zeit \texorpdfstring{$\tau(K_6)$}{\tau(K_6)}}

Zeit entsteht aus der \emph{Stabilität kognitiver Zyklen}:
\[
\tau(K_6)
= \min_j \tau_{C_j},
\qquad
\Pi(C_j) > \Theta_{\mathrm{Zeit}}.
\]

Die kognitive Zeit ist langsamer und integrativer als die elektrische Zeit ($K_5$).

% ================================================================
\subsubsection{Kontinuität \texorpdfstring{$k(K_6)$}{k(K_6)}}

Unter Verwendung der allgemeinen Formel (K6 Lauf 7):

\[
k_6
= H(\Omega(K_6))
  \cdot \frac{|C_{\max}^{(6)}|}{|C_{\mathrm{all}}|}
  \cdot \frac{r}{r_{\max}}
  \cdot \left(1 - \frac{D(m)}{\Theta_{\mathrm{cog,max}}}\right)_+
  \cdot \left(1-\frac{\sigma(J)}{\sigma_{\max}}\right).
\]

Wo:
\begin{itemize}
    \item $|C_{\max}^{(6)}|$ — Größe der größten verbundenen Komponente in $G_{\mathrm{cog}}$,
    \item $r$ – Anzahl der aktiven Achsen vs. maximale Achsen,
    \item $D(m)$ – Modellinkonsistenzmetrik.
\end{itemize}

Interpretation:
\begin{itemize}
    \item große verbundene semantische Struktur erhöht $k_6$,
    \item-Inkonsistenz und Vorhersagefehler reduzieren $k_6$,
    \item Fragmentierung des kognitiven Diagramms kollabiert $k_6$.
\end{itemize}

% ================================================================
\subsubsection{Strukturelle Spannung \texorpdfstring{$T(K_6)$}{T(K_6)}}

Quellen:
\begin{itemize}
    \item hoher Vorhersagefehler $PE_t$,
    \item inkompatible Bindungen,
    \item semantische Konflikte,
    \item instabile Speicherneukonsolidierung,
    \item inkohärenter Modellwettbewerb,
    \item außer Kontrolle geratene Salienz.
\end{itemize}

Zusammenbrechen, wenn:
\[
T(K_6) > \Theta_{\mathrm{cog}}.
\]

% ================================================================
\subsubsection{Energie \texorpdfstring{$E(K_6)$}{E(K_6)}}

Energie ist ein kognitiv-neuraler Hybrid:
\[
E(K_6)
= E_{K_5}
+ E_{\mathrm{pred}}
+ E_{\mathrm{bind}}
+ E_{\mathrm{mem}}
+ E_{\mathrm{cmp}}
+ E_{\mathrm{Modell}}.
\]

Interpretation:
\begin{itemize}
    \item Vorhersage-Aktualisierungsenergie,
    \item Speichererhaltungsenergie,
    \item verbindliche Energiekosten,
    \item Energie in repräsentativen Diagrammen,
    \item Kosten für den Modellwechsel.
\end{itemize}

% ================================================================
\subsubsection{Operatoren auf \texorpdfstring{$K_6$}{K_6} (\texorpdfstring{$\Psi$}{\Psi}, \texorpdfstring{$\Phi$}{\Phi}, \texorpdfstring{$\Lambda$}{\Lambda}, \texorpdfstring{$U$}{U}, \texorpdfstring{$\Chi$}{\Chi})}

\begin{itemize}
    \item $\Psi_{6\to7}$ – Entstehung kommunikativer Achsen (gemeinsame Modelle),
    \item $\Phi$ – Entwicklung der Repräsentationsstruktur,
    \item $\Lambda$ – Vereinheitlichung von Modellkomponenten in Strukturen höherer Ordnung,
    \item $U$ – Integration kognitiver Zyklen in stabile semantische Einheiten,
    \item $\Chi$ – Verzweigung in Spezialisierungen (speicherlastig, vorhersagelastig, bindungslastig).
\end{itemize}

% ================================================================
\subsubsection{Prozesse auf \texorpdfstring{$K_6$}{K_6}}

Typisch:
\begin{itemize}
    \item Modellbildung und -auflösung,
    \item Vorhersage und Fehlerkorrektur,
    \item semantische Bindung,
    \item Konsolidierung und Rückkonsolidierung,
    \item Salienzmodulation,
    \item Merkmalsextraktion,
    \item Generierung von Proto-Symbolen $\sigma$,
    \item Attraktordynamik prägt die kognitive Topologie.
\end{itemize}

% ================================================================
\subsubsection{Vorhersagen für \texorpdfstring{$K_6$}{K_6}}

\begin{enumerate}
    \item Ein kognitiver Zusammenbruch tritt auf, wenn sich $PE_t$ schneller ansammelt, als sich die Bindung stabilisieren kann.
    \item Ein stabiles Langzeitgedächtnis erfordert Rekonsolidierungszyklen.
    \item Bindungsachsen haben eine maximale Kapazität, die experimentell getestet werden kann.
    \item Semantische Konnektivität sagt Robustheit gegenüber Rauschen voraus.
    \item Die kognitive Zeit $\tau(K_6)$ nimmt mit der Komplexität der Darstellung zu.
    \item Protosymbolische Zustände treten auf natürliche Weise auf, wenn $G_{\mathrm{cog}}$ eine ausreichende Konnektivität erreicht.
\end{enumerate}

% ================================================================
\subsubsection{Experimente für \texorpdfstring{$K_6$}{K_6}}

Mögliche empirische Analoga:
\begin{itemize}
    \item neuronale Netzwerk-Attraktor-Modelle mit kontrolliertem Rauschen,
    \item prädiktive Codierungsexperimente zur Messung von $PE_t$-Schwellenwerten,
    \item Bindungskapazitätsverhaltenstests,
    \item Protokolle zur Speicherrekonsolidierung,
    \item semantische Graphfragmentierungstests,
    \item untersucht die Entstehung von Protosymbolen in neuronalen Systemen.
\end{itemize}

% ================================================================
\subsubsection{Zusammenbruch und Tod von \texorpdfstring{$K_6$}{K_6}}

Fehler, wenn:
\[
\Omega(K_6) = \varnothing.
\]

Mechanismen:
\begin{itemize}
    \item unkontrolliertes Wachstum des Vorhersagefehlers,
    \item Katastrophaler Bindungsfehler,
    \item Zusammenbruch des semantischen Diagramms,
    \item Gedächtniszerfall,
    \item außer Kontrolle geratene Salienzkaskaden,
    \item Unfähigkeit, Modellkonflikte zu lösen.
\end{itemize}

Der Tod von $K_6$ verhindert den Übergang zur sozialen Kognition $K_7$.

% ================================================================
\subsubsection{Falschbarkeit von \texorpdfstring{$K_6$}{K_6}}

Das Modell sagt Folgendes voraus:
\begin{itemize}
    \item Existenz strenger Vorhersagefehlerschwellen,
    \item Notwendigkeit von Bindungskapazitätsgrenzen,
    \item messbare kognitive Zyklen (Vorhersage, Gedächtnis, Bindung),
    \item hierarchische Entstehung von Protosymbolen,
    \item Abhängigkeit der semantischen Stabilität von der Graphkonnektivität.
\end{itemize}

Eine Widerlegung würde sich aus der Beobachtung kognitiver Systeme ohne diese Strukturen ergeben.

% ================================================================
\subsubsection{Verzweigung / ontologische Position von \texorpdfstring{$K_6$}{K_6}}

Filialen:
\begin{itemize}
    \item speicherdominante vs. vorhersagedominante Systeme,
    \item hochbindende vs. niedrigbindende Organismen,
    \item modellreiche vs. modellarme kognitive Architekturen,
    \item verteilte vs. lokalisierte semantische Graphen.
\end{itemize}

Ontologische Positionierung:
\[
K_5 \to K_6 \to K_7.
\]

$K_6$ ist die Voraussetzung für Kommunikation, Sozialverhalten und gemeinsame Modelle.

% ================================================================
\subsubsection{Beziehung zu M-Räumen}

$M_6$ Einschränkungen:
\begin{itemize}
    \item zulässige Vorhersagefehlerlandschaften,
    \item Bindungstopologiebereiche,
    \item Speicherstabilitätsregime,
    \item strukturelle Kohärenzanforderungen,
    \item zulässige Dynamik repräsentativer Graphen.
\end{itemize}

Die Kompatibilität mit $M_6$ bestimmt die Lebensfähigkeit kognitiver Kontinua.
